\chapter*{前言}
\markboth{前言}{前言}
\addcontentsline{toc}{chapter}{前言}
\label{chap:preface}

\SourceParagraph{preface-001}{7}
2002年 Clay 几何与弦理论讲习班于2002年3月25日至4月19日在英国剑桥 的 Isaac Newton 数学科学研究所 举行。讲习班由当时在 Newton 研究所 同期举办的两个研讨会的组织者共同主持：其中，高维复几何研讨会由 Alessio Corti、Mark Gross 和 Miles Reid 组织，M 理论研讨会由 Robbert Dijkgraaf、Michael R. Douglas、Jerome Gauntlett 和 Chris Hull 组织；时任 Clay 数学研究所 主席的 Arthur Jaffe 也参与了合作。

\SourceParagraph{preface-002}{7}
记录这次讲习班学术成果的书籍共有两卷，本书是其中的第二卷。第一卷是由 Michael R. Douglas、Jerome Gauntlett 和 Mark Gross 编辑的《弦与几何》\BookCite{ref131}，属于会议论文集，主要讨论特殊完整群流形和超引力。

\SourceParagraph{preface-003}{7}
本卷旨在成为一部专著，从同调和环面纤维化两个观点介绍镜像对称。我们希望，本卷能自然地接续 Hori、Katz、Klemm、Pandharipande、Thomas、Vafa、Vakil 和 Zaslow 合著的《镜像对称》\BookCite{ref242}；那部书是2000年春季第一届 Clay 讲习班的成果。以下我们将其简称为 \BookBibText{ref242}{MS1}。熟悉 \BookBibText{ref242}{MS1} 中的基础内容，可以视为阅读本卷的先决条件；我们也将经常引用 \BookBibText{ref242}{MS1}，供读者查阅背景知识。

\SourceParagraph{preface-004}{7}
本卷的总体目标是探讨 Dirichlet 膜的物理和数学两个方面。全书的叙述围绕两个主要思想展开：Kontsevich 的同调镜像对称猜想和 Strominger--Yau--Zaslow 猜想。虽然 Kontsevich 提出这一猜想时，D 膜还没有被引入物理学，但我们将说明，这一猜想实际上等价于将两种不同的 D 膜范畴等同起来。特别地，我们将考察物理如何自然地引导我们走向导出范畴、Fukaya 范畴等数学概念。我们会探究 Strominger--Yau--Zaslow 猜想的各种影响及其研究现状，并将这些思想与若干活跃的研究领域联系起来，包括 McKay 对应、拓扑量子场论和稳定性结构。

\SourceParagraph{preface-005}{8}
同整个镜像对称领域一样，这些领域也从数学家与物理学家之间成果丰硕的互动中获益良多。在本书历时六年的酝酿过程中，人们在澄清和理解最初的猜想方面取得了大量进展，并且在某些情形下证明了这些猜想。

\SourceParagraph{preface-006}{8}
可以说，要充分领会由此形成的图景，就需要对数学和物理两个方面的观点都有一定了解。在同一本书中传达这两种观点，既是一项挑战，也是一个机会。我们并不满足于把同一件事讲述两遍，分别采用一种观点；相反，我们尝试给出统一的阐述，使数学和物理各自贡献其不可或缺的见解，同时让物理学家和数学家即使未必熟练掌握这两个学科的全部基础，也能读懂。

\SourceParagraph{preface-007}{8}
完成这一任务的部分困难，来自物理学家与数学家在背景知识和语言上的诸多差异。虽然我们觉得自己已经为弥合这些差异做了大量工作，但也因此更加清楚地看到了仍然存在的许多隔阂。

\SourceParagraph{preface-008}{8}
当然，还有一个更根本的困难：要讲述完整的故事，所需涉及的主题十分广泛，以至于没有一位作者是所有这些领域的专家。因此，我们将主要撰写工作分工进行，同时仍力求通过大量编辑和交叉引用，使全书成为一个统一的整体。这项工作由 Michael R. Douglas 和 Mark Gross 根据所有作者相互审读后提出的意见完成。本书在这一层面上的成败，由 Michael R. Douglas 和 Mark Gross 承担责任。

\SourceParagraph{preface-009}{8}
\BookChapterRef{1}旨在主要从物理的角度概览全书主题。\BookChapterRef{2}讨论拓扑开弦理论，由 Greg Moore 和 Graeme Segal 撰写。这部分内容的早期稿件发表于 \href{https://arxiv.org/abs/hep-th/0609042v1}{arXiv:hep-th/0609042v1}。

\SourceParagraph{preface-010}{8}
\BookChapterRef{3}和\BookChapterRef{5}讨论 Dirichlet 膜的物理，主要由 Paul Aspinwall、Michael R. Douglas 和 Anton Kapustin 撰写。部分内容曾发表于 \href{https://arxiv.org/abs/hep-th/0403166}{arXiv:hep-th/0403166}，而\BookSectionRef{sec:5.7}{5.7}和\BookSectionRef{sec:5.8}{5.8}则在很大程度上以 Tom Bridgeland 已发表的稳定性结构研究为基础。

\SourceParagraph{preface-011}{8}
\BookChapterRef{4}讨论表示论，主要由 Tom Bridgeland、Alastair Craw 和 Balázs Szendrői 撰写。

\SourceParagraph{preface-012}{8}
\BookChapterRef{6}和\BookChapterRef{8}由 Mark Gross 撰写，\BookChapterRef{7}由 Mark Gross 和 Pelham Wilson 撰写。

\SourceParagraph{preface-013}{8}
Robert Karp 和 Arthur Greenspoon 阅读了全部书稿，二人都发现了许多不够精确和不够清晰之处。我们也从与 Mohammed Abouzaid、Gary Gibbons、Akira Ishii、Dmitri Orlov 和 Bernd Siebert 的讨论及他们的意见中获益。几位作者还希望感谢 IHES 的热情接待，本书的部分内容是在那里完成的。

\SourceParagraph{preface-014}{9}
让我们再次向促成2002年讲习班举办的诸位表示感谢：Newton 研究所 的两位主任 H. Keith Moffatt 和 John Kingman，以及该所工作人员 Wendy Abbott、Tracey Andrew、Caroline Fallon、Jackie Gleeson、Louise Grainger、Rebecca Speechley 和 Christine West。

\SourceParagraph{preface-015}{9}
最后，我们感谢 Jim Carlson 以及 Clay 数学研究所 的工作人员，尤其是 Vida Salahi，在本书的筹备和制作过程中提供的帮助。我们特别感谢 Lori Lejeune 为本书提供插图。

\SourceParagraph{preface-016}{9}
\begin{flushright}
Michael R. Douglas 和 Mark Gross\\
2008年10月
\end{flushright}
